\section{Correlated State Variables}

As an alternative to using all technical indicators, the 11 indicators that correlate most
with tomorrow's log-return of the closing price were selected as state variables. This
reduces the dimensionality of the state space while aiming to retain the most informative
features.

\begin{figure}[H]
  \centering
  \includegraphics[width=0.78\textwidth]{z122_1.png}
  \caption{Correlated state variables -- cumulative return ratio: buy and hold vs.\ agent.}
\end{figure}

\begin{figure}[H]
  \centering
  \includegraphics[width=0.78\textwidth]{z123_1.png}
  \caption{Correlated state variables -- actions chosen by the trained agents over time.}
\end{figure}

\paragraph{Comments.}
Comparing the graphs on the ratio of revenues between buy and hold and the trained agent,
it is obvious that the revenues are clearly better for the trained agent in comparison to
buy and hold. However, the actions chosen over time give an insight into the fact that the
trained agents suggest clearly different decisions. It can therefore be assumed that the
trained agent applied to this share could lead to inconsistent behavior. Comparing those
results with the agent trained with all indicators from above, it is evident that in the
case of positive revenues, the trained agent using the 11 indicators that correlate most
with tomorrow's log-return of the close price provides similar results as the one using
all indicators.
